\documentclass{beamer}
\usetheme{Madrid}
\usecolortheme{default}

\usepackage[T1]{fontenc}
\usepackage[utf8]{inputenc}
\usepackage{booktabs}

\title{RNNs and Backpropagation}
\subtitle{How This Model Reads a Sequence and Learns from Loss}
\author{}
\date{}

\begin{document}

\frame{\titlepage}

\begin{frame}{RNN Process}
\small
The RNN path in this project is:

\vspace{0.2cm}

\begin{center}
\texttt{image sequence -> RNN loop -> final memory -> digit scores}
\end{center}

\vspace{0.25cm}

\begin{itemize}
    \item The image is treated like a sequence of input steps.
    \item The RNN reads one step at a time.
    \item It keeps a memory called the hidden state.
    \item After the last step, that memory is used to predict the digit.
\end{itemize}
\end{frame}

\begin{frame}[fragile]{Step 1: Input Shape}
\small
In \texttt{model/model.py}, the model starts with:

\vspace{0.15cm}

\begin{verbatim}
B, L, C = inputs.shape
\end{verbatim}

\vspace{0.2cm}

\begin{center}
\begin{tabular}{l l}
\toprule
Name & Meaning \\
\midrule
\texttt{B} & batch size: how many images \\
\texttt{L} & sequence length: how many steps \\
\texttt{C} & channels: numbers per step \\
\bottomrule
\end{tabular}
\end{center}

For grayscale MNIST, each pixel step usually has one brightness value, so \texttt{C = 1}.
\end{frame}

\begin{frame}[fragile]{Step 2: Put Sequence First}
\small
The input starts as:

\vspace{0.15cm}

\begin{center}
\texttt{[batch, length, channels]}
\end{center}

\vspace{0.15cm}

Then the code changes it:

\begin{verbatim}
inputs = inputs.transpose(0, 1)
\end{verbatim}

\vspace{0.15cm}

Now it is:

\begin{center}
\texttt{[length, batch, channels]}
\end{center}

This makes it easy for the RNN to read step 1, then step 2, then step 3.
\end{frame}

\begin{frame}[fragile]{Step 3: Start with Empty Memory}
\small
Before reading the sequence, the RNN creates the first hidden state:

\vspace{0.15cm}

\begin{verbatim}
state = self.cell.default_state(inputs[0], max_batch_size)
\end{verbatim}

\vspace{0.2cm}

\begin{center}
\texttt{state = empty memory}
\end{center}

\vspace{0.2cm}

At the beginning, the RNN has not seen any pixels yet, so its memory starts as zeros.
\end{frame}

\begin{frame}[fragile]{Step 4: Loop Through the Sequence}
\small
In \texttt{model/rnn.py}, the RNN reads one input step at a time:

\vspace{0.15cm}

\begin{verbatim}
for input in torch.unbind(inputs, dim=0):
    output, new_state = self.cell.forward(input, state)
    state = new_state
\end{verbatim}

\vspace{0.2cm}

\begin{center}
\texttt{current step + old memory -> new memory}
\end{center}

\vspace{0.15cm}

This is the main RNN idea: the next step depends on what the model already remembers.
\end{frame}

\begin{frame}[fragile]{Step 5: Update the Memory}
\small
The actual memory update happens in \texttt{model/rnncell.py}:

\vspace{0.15cm}

\begin{verbatim}
hidden_preact = self.W_hx(input) + self.W_hh(h)
hidden = self.hidden_activation_fn(hidden_preact)

return hidden, hidden
\end{verbatim}

\vspace{0.15cm}

\begin{itemize}
    \item \texttt{self.W\_hx(input)} looks at the current input step.
    \item \texttt{self.W\_hh(h)} looks at the previous memory.
    \item The activation turns the combined value into the new memory.
\end{itemize}
\end{frame}

\begin{frame}{Step 6: What the Activation Does}
\small
The activation is applied after input information and memory information are added.

\vspace{0.2cm}

\begin{center}
\texttt{new memory = activation(input part + memory part)}
\end{center}

\vspace{0.2cm}

\begin{itemize}
    \item It turns raw numbers into hidden-state values.
    \item It makes the model nonlinear, so it can learn shapes and patterns.
    \item It controls how strongly each memory feature turns on or off.
\end{itemize}
\end{frame}

\begin{frame}{Why \texttt{tanh}?}
\small
Vanilla RNNs often use \texttt{tanh} for the hidden activation.

\vspace{0.2cm}

\begin{center}
\texttt{large negative -> near -1}

\texttt{0 -> 0}

\texttt{large positive -> near 1}
\end{center}

\vspace{0.2cm}

\begin{itemize}
    \item It keeps memory values between \texttt{-1} and \texttt{1}.
    \item It is centered around zero.
    \item It allows negative, positive, and neutral memory values.
\end{itemize}
\end{frame}

\begin{frame}[fragile]{Step 7: Use the Final Memory}
\small
After all input steps are read, \texttt{model/model.py} keeps the final state:

\vspace{0.15cm}

\begin{verbatim}
_, state = self.rnn(inputs, init_state=initial_state,
                    return_output=False)
state = self.rnn.output(state)
return self.output_mlp(state)
\end{verbatim}

\vspace{0.2cm}

\begin{center}
\texttt{final memory -> output layer -> 10 digit scores}
\end{center}

So the RNN summarizes the image first, then the classifier predicts the digit.
\end{frame}

\begin{frame}{From Pixels to Features}
\small
The model does not jump directly from pixels to the answer.

\vspace{0.25cm}

\begin{center}
\texttt{784 pixel brightness values -> 256 features -> 10 scores}
\end{center}

\vspace{0.25cm}

\begin{itemize}
    \item The 256 features are learned numbers that summarize the image.
    \item \texttt{256} comes from \texttt{hidden\_size: 256}.
    \item If \texttt{hidden\_size} were \texttt{64}, the model would use 64 features.
\end{itemize}
\end{frame}

\begin{frame}[fragile]{Final Output Layer}
\small
In \texttt{model/model.py}, the final output layer is built here:

\vspace{0.15cm}

\begin{verbatim}
backbone_output_size = self.rnn.output_size()
sizes = [backbone_output_size] + output_hiddens + [output_size]
self.output_mlp = nn.Sequential(
    *[nn.Linear(sizes[i], sizes[i+1])
      for i in range(len(sizes)-1)]
)
\end{verbatim}

\vspace{0.15cm}

For MNIST with \texttt{hidden\_size = 256}, this becomes:

\begin{center}
\texttt{nn.Linear(256, 10)}
\end{center}
\end{frame}

\begin{frame}{Linear Layer Meaning}
\small
A linear layer combines features with learned weights:

\[
\text{scores} = \text{features} \times W + b
\]

\vspace{0.15cm}

\begin{center}
\begin{tabular}{l l}
\toprule
Part & Shape \\
\midrule
\texttt{features} & \texttt{(256)} \\
\texttt{W} conceptually & \texttt{(256, 10)} \\
\texttt{b} & \texttt{(10)} \\
\texttt{scores} & \texttt{(10)} \\
\bottomrule
\end{tabular}
\end{center}

PyTorch stores the linear weight as \texttt{(10, 256)}, but the idea is the same.
\end{frame}

\begin{frame}[fragile]{Batch, X, Y, and Training Step}
\small
\begin{center}
\begin{tabular}{l l}
\toprule
Name & Meaning \\
\midrule
\texttt{batch\_x} & input images as pixel sequences \\
\texttt{batch\_y} & correct digit labels \\
\texttt{out} & model output scores \\
\texttt{loss} & how wrong the scores are \\
\bottomrule
\end{tabular}
\end{center}

\vspace{0.15cm}

\begin{verbatim}
batch_x, batch_y, *len_batch = batch
out = self.model(batch_x, len_batch)
loss = self.dataset.loss(out, batch_y, len_batch)
return loss
\end{verbatim}
\end{frame}

\begin{frame}{Forward Pass and Output Scores}
\small
The forward pass means the model makes a prediction using its current weights.

\vspace{0.25cm}

\begin{center}
\texttt{pixels -> model -> 10 digit scores}
\end{center}

\vspace{0.25cm}

Suppose the image shows the digit \texttt{3}:

\vspace{0.15cm}

\begin{center}
\begin{tabular}{c c c c c c c c c c}
0 & 1 & 2 & 3 & 4 & 5 & 6 & 7 & 8 & 9 \\
\midrule
0.1 & 0.2 & -0.5 & 5.0 & 0.0 & -1.0 & 0.3 & 0.4 & -0.7 & 0.2
\end{tabular}
\end{center}

The highest score is for digit \texttt{3}, so the model predicts \texttt{3}.
\end{frame}

\begin{frame}{Weights for Digit 3}
\small
Each digit score is a weighted combination of the 256 features.

\[
\text{score for 3}
=
f_1w_1 + f_2w_2 + \cdots + f_{256}w_{256} + b
\]

\vspace{0.2cm}

\begin{itemize}
    \item The weights connected to \texttt{score\_3} are the weights for digit \texttt{3}.
    \item At first, these weights are random.
    \item During training, they become useful for making images of \texttt{3} produce high \texttt{score\_3}.
\end{itemize}
\end{frame}

\begin{frame}[fragile]{How Digit 3 Is Learned}
\small
Suppose the image is actually a \texttt{3}:

\begin{center}
\texttt{batch\_x = image pixels}, \qquad \texttt{batch\_y = 3}
\end{center}

If the model predicts \texttt{7}, then \texttt{score\_7} is too high,
\texttt{score\_3} is too low, and the loss becomes high.

\vspace{0.15cm}

\begin{verbatim}
out = self.model(batch_x, len_batch)
loss = self.dataset.loss(out, batch_y, len_batch)
return loss
\end{verbatim}

Backpropagation uses this loss to change the weights a little.
\end{frame}

\begin{frame}{Repeated Learning}
\small
One example changes the weights only a little.

\vspace{0.25cm}

\begin{center}
\texttt{pixels -> scores -> loss -> gradients -> update}
\end{center}

\vspace{0.25cm}

\begin{itemize}
    \item After many images of \texttt{3}, \texttt{score\_3} becomes better for digit \texttt{3}.
    \item After many images of other digits, the model learns when not to make \texttt{score\_3} high.
    \item Learning happens over many batches and epochs.
\end{itemize}
\end{frame}

\begin{frame}[fragile]{Loss}
\small
Loss is one number that says how wrong the model is.

\begin{itemize}
    \item Low loss means the prediction is good.
    \item High loss means the prediction is bad.
    \item For MNIST, this project uses cross entropy loss.
\end{itemize}

\vspace{0.15cm}

\begin{verbatim}
class MulticlassClassification(Task):
    @staticmethod
    def loss(logits, y, len_batch=None):
        return F.cross_entropy(logits, y)
\end{verbatim}
\end{frame}

\begin{frame}{Chain Rule and Backpropagation}
\small
Backpropagation uses the chain rule to calculate gradients.

\[
\frac{dL}{dx}
=
\frac{dL}{dy}
\times
\frac{dy}{dx}
\]

\vspace{0.15cm}

For a model:

\[
\text{weight} \rightarrow \text{feature} \rightarrow \text{score} \rightarrow \text{loss}
\]

\[
\frac{dLoss}{dWeight}
=
\frac{dLoss}{dScore}
\times
\frac{dScore}{dFeature}
\times
\frac{dFeature}{dWeight}
\]
\end{frame}

\begin{frame}{Backward Pass}
\small
The backward pass starts from the loss and moves backward through the model.

\vspace{0.25cm}

\begin{center}
\texttt{loss -> scores -> features -> weights}
\end{center}

\vspace{0.25cm}

\begin{itemize}
    \item It uses the chain rule.
    \item It calculates gradients for the weights.
    \item This is backpropagation.
\end{itemize}
\end{frame}

\begin{frame}{Weight Update}
\large
\[
\text{new weight}
=
\text{old weight}
-
\text{learning rate}
\times
\text{gradient}
\]

\vspace{0.35cm}

\normalsize
\begin{itemize}
    \item Backpropagation calculates gradients.
    \item The optimizer uses gradients to update weights.
    \item The new weight can become smaller or larger.
    \item The goal is smaller loss, not smaller weights.
\end{itemize}
\end{frame}

\end{document}
